% ===== 原文 pp.142--149（PDF p.152--159）：附录 C 相伴纤维丛与等变上同调 =====
% ===== 本文件 appC-01.tex：印刷 pp.142--145（PDF p.152--155） =====
% 本附录无编号公式、无图；全部显示公式用 \[ … \]，一律不编号。

% [原文 p.142]
\chapter{相伴纤维丛与等变上同调}

规范 $\sigma$-模型在本论文中起着重要作用：带规范场的 4d A-模型在第 3 章讨论，带规范场的 3d A-模型与带规范场的 Rozansky-Witten 模型在第 5 章讨论，而修正的 2d A-模型与带规范场的 2d B-模型在第 7 章讨论。

由于相伴纤维丛的概念为规范 $\sigma$-模型提供了几何框架，我们在这里回顾相伴纤维丛的数学，特别是解释如何用靶空间上的等变上同调类来写出规范 $\sigma$-模型作用量中的拓扑项。定理大多不加证明地陈述。

% [Tutor Note: p.142 术语首现：associated fiber bundle = 相伴纤维丛（glossary 已收）；gauged sigma model = 规范 $\sigma$-模型。后文出现的 “basic form”（基本形式）指同时不变且水平的微分形式，是 Cartan 模型下拉回能下降到基流形的关键性质（见 p.148 性质 1）。]

设 $P$ 是底流形 $M$ 上的主 $G$-丛，其右 $G$-作用为 $r: P\times G\to P$；又设 $X$ 是配有左 $G$-作用的流形
\[ l: G\times X\to X. \]
我们用点 $\cdot$ 同时表示 $l$ 与 $r$ 两种作用。在物理应用中，$M$ 是时空，$P$ 是 $M$ 上的规范丛，而 $X$ 是 $\sigma$-模型的靶空间。五元组 $(M,P,G,l,X)$ 确定 $M$ 上一个以 $X$ 为典型纤维、以 $G$ 为结构群的纤维丛，称为\emph{相伴纤维丛}。我们依次讨论关于这个相伴纤维丛“是什么”的两种互补观点：全空间观点与主丛观点。

\section{相伴纤维丛的全空间}

我们假定 $M$ 是 $2m$ 维的 Kähler 流形，$X$ 是 $2n$ 维的 Kähler 流形，并且 $X$ 上的 $G$-作用保持其度规 $g_X$、复结构 $J_X$ 与 Kähler 形式 $\Omega_X$。我们后面要说的大部分内容在 $M$ 与 $X$ 没有 Kähler 结构与复结构时也成立；这些结构只在全空间上构造 Kähler 结构与复结构时才是必需的。下面我们描述相伴纤维丛的全空间与投影映射。

% [原文 p.143]
\begin{definition}
相伴纤维丛 $(M,P,G,l,X)$ 的全空间 $E$ 是 $(2m+2n)$ 维流形 $E=P\times_G X$，它定义为商空间
\[ E = P\times_G X = (P\times X)/G, \]
其中 $G$ 在 $P\times X$ 上的右作用由 $(p,x)\cdot g\equiv(p\cdot g,\ g^{-1}\cdot x)$ 给出，对一切 $p\in P$、$x\in X$ 与 $g\in G$。
\end{definition}

\begin{definition}
（满的）投影映射 $\pi_E: E\to M$ 是由
\[ \pi_E[p,x] = \pi(p) \]
给出的映射，对一切 $p\in P$ 与 $x\in X$，其中 $\pi: P\to M$ 是底层主丛的投影映射。
\end{definition}

% [Tutor Note: p.143 商空间 $E=(P\times X)/G$ 是光滑流形、且 $\dim E = \dim P+\dim X = 2m+2n$，这两点依赖右作用的自由性：主丛的 $G$-作用在 $P$ 上自由，故 $(p,x)\cdot g=(p,x)$ 迫使 $g=e$，于是 $G$ 在 $P\times X$ 上的作用也自由，商才是流形而不会出现轨形点。核对：取 $X$ 为一点，$(P\times X)/G = P/G = M$，即 $E$ 退化为底流形。]

把切空间 $TE$ 明确描述出来也会有用：

\begin{theorem}
$E$ 的切丛的全空间 $TE$ 简单地由扭积
\[ TE = TP\times_G TX \]
给出，它由切向量的等价类
\[ TE = \{\,[u_p,w_x]\ \mid\ u_p\in T_pP,\ w_x\in T_xX\,\}, \]
构成，其中
\[
    [u_p,w_x] = [\dd r_g u_p,\ \dd l_{g^{-1}}w_x] = [u'_{p\cdot g},\ w'_{g^{-1}\cdot x}].
\]
切向量 $u'_{p\cdot g}\in T_{p\cdot g}P$ 与 $w'_{g^{-1}\cdot x}\in T_{g^{-1}\cdot x}X$ 分别是 $u_p$ 与 $w_x$ 在 $P$ 上右作用 $g$、在 $X$ 上左作用 $g^{-1}$ 之下的前推。此外，$TE$ 配有显然的投影 $TE\to E$，它把 $[u_p,w_x]$ 送到 $[p,x]$；这是 $E$ 上的向量丛，每个纤维 $T_{[p,x]}E$ 都是 $2m+2n$ 维向量空间。
\end{theorem}

% [Tutor Note: p.143 该定理的核对办法：$TE$ 的纤维维数应为 $\dim T_pP+\dim T_xX = 2m+2n$，与 $E$ 的维数相符；且 $[u_p,w_x]$ 的等价关系必须与 $T_pP\to T_{p\cdot g}P$、$T_xX\to T_{g^{-1}\cdot x}X$ 两处推前一致，这正是扭积 $TP\times_G TX$ 的定义。自洽检验：当 $X$ 为一点时 $TX=0$，应退化为 $TP\times_G 0 = TM$，与“$E=M$ 时 $TE=TM$”相符。]

\begin{definition}
令 $E^{(\alpha)}\equiv\pi_E^{-1}(U^{(\alpha)})$ 与 $E_\sigma\equiv\pi_E^{-1}(\sigma)$，其中 $U^{(\alpha)}\subset M$ 是 $P$ 的一个平凡化邻域，$\sigma\in U^{(\alpha)}$ 是该邻域中的一点。也就是说，$E^{(\alpha)}$ 是 $E$ 位于 $U^{(\alpha)}$ 上方的部分，而 $E_\sigma$ 是 $\sigma$ 上方的纤维。
\end{definition}

\begin{definition}
令 $x^{(\alpha)}: E^{(\alpha)}\to X$ 为丛 $E$ 的平凡化映射，它以丛 $P$ 的平凡化映射 $g^{(\alpha)}: P^{(\alpha)}\to G$ 定义如下：
\[ x^{(\alpha)}[p,x] = (g^{(\alpha)}(p))^{-1}\cdot x . \]
令 $x^{(\alpha)}_\sigma: E_\sigma\to X$ 为 $x^{(\alpha)}$ 到 $E_\sigma\subset E^{(\alpha)}$ 的限制。
\end{definition}

% [原文 p.144]
映射 $x^{(\alpha)}$ 给出 $E$ 在 $U^{(\alpha)}$ 上的一个平凡化：
\[ E^{(\alpha)}\simeq U^{(\alpha)}\times X, \]
经由映射 $\simeq: [p,x]\mapsto(\pi_E[p,x],\ x^{(\alpha)}[p,x]) = (\pi(p),\ (g^{(\alpha)}(p))^{-1}\cdot x)$。限制映射
\[ x^{(\alpha)}_\sigma \equiv x^{(\alpha)}|_{E_\sigma}: E_\sigma\to X \]
是微分同胚；也就是说，每个纤维 $E_\sigma$ 在微分同胚 $x^{(\alpha)}_\sigma$ 之下微分同胚于 $X$ 的一个副本。正如标号 $(\alpha)$ 所表明的，这个具体的微分同胚是由底层主丛 $P$ 的平凡化诱导的。

% [Tutor Note: p.144 平凡化映射之所以必须写成 $(g^{(\alpha)}(p))^{-1}\cdot x$（左作用，且取逆）：$E$ 是左作用与右作用的混合商，只有这样才能保证 $x^{(\alpha)}[p\cdot h, h^{-1}\cdot x]=x^{(\alpha)}[p,x]$，即良定义。核对：$g^{(\alpha)}(p\cdot h)=g^{(\alpha)}(p)h$ 时，$(g^{(\alpha)}(p)h)^{-1}\cdot(h^{-1}\cdot x)=(g^{(\alpha)}(p))^{-1}\cdot x$。]

仿照 $P$ 的垂直向量的概念，我们定义

\begin{definition}
垂直子丛 $VE\subset TE$ 由那些被投影映射的前推映为零的切向量 $[u_p,w_x]\in T_{[p,x]}E$ 组成：
\[ VE \equiv \ker \dd\pi_E . \]
空间 $VE$ 是 $E$ 上的向量丛，其典型纤维为 $V_{[p,x]}E$，一个 $2n$ 维向量空间；也就是说，$V_{[p,x]}E$ 由沿纤维 $E_\sigma$ 的切向量组成，其中 $\sigma=\pi_E[p,x]=\pi(p)$：
\[ V_{[p,x]}E = T_{[p,x]}(E_{\pi(p)})\subset T_{[p,x]}E . \]
\end{definition}

% [Tutor Note: p.144 维数核对：$\dim V_{[p,x]}E = \dim T_{[p,x]}E - \dim T_{\pi(p)}M = (2m+2n)-2m = 2n$，与“纤维是 $X$ 的副本、$\dim X=2n$”一致，这是 $VE$ 定义正确的自洽检验。]

我们可以用微分同胚 $x^{(\alpha)}_\sigma$ 把 $X$ 上的度规、复结构与 Kähler 形式拉回到 $VE$ 的各纤维上，得到相应的度规 $g^V$、复结构 $J^V$ 与辛 2-形式 $\Omega^V$：

\begin{definition}
张量
\[
    g^V|_e\in S^2(V_eE),\qquad J^V|_e\in V_e^*E\otimes V_eE,\qquad \Omega^V|_e\in\Lambda^2(V_eE)
\]
由 $\dd x^{(\alpha)}_\sigma: V_eE\to T_{x^{(\alpha)}_\sigma(e)}X$ 之下的拉回给出，其中 $\sigma=\pi_E(e)$ 而 $e=[p,x]\in E$。我们无须在 $g^V|_e$ 上加标号 $(\alpha)$，正是因为在假定之下 $g_X$ 是 $G$-不变的。
\end{definition}

% [Tutor Note: p.144 “无须加标号 $(\alpha)$”这一句的核对：换用另一个平凡化 $(\beta)$ 时两个微分同胚 $x^{(\alpha)}_\sigma$ 与 $x^{(\beta)}_\sigma$ 相差一个左 $G$-作用，而 $g_X$ 的 $G$-不变性使两者拉回的度规相同；同理 $J_X$、$\Omega_X$ 的 $G$-不变性保证 $J^V$、$\Omega^V$ 也无标号。这也解释了 §C.1 开头为何要假定 $G$ 保持 $g_X$、$J_X$、$\Omega_X$。]

至此我们有了度规、（殆）复结构与 Kähler 形式的概念，但它们严格地只作用在垂直向量上。要把这些结构的作用延拓到 $TE$ 的所有方向，我们需要在底层主丛 $P$ 上选取一个联络 $\omega$。这样的联络给出 $P$ 上垂直投影与水平投影的概念，进而给出 $E$ 上垂直投影与水平投影的概念：

% [原文 p.145]
\begin{definition}
令 $\omega^V_E: TE\to VE$ 为由
\[ \omega^V_E[u_p,w_x]\equiv[\omega^V u_p,\ w_x]\in V_{[p,x]}E \]
给出的垂直投影，其中 $u_p\in T_pP$、$w_x\in T_xX$，而 $\omega^V: TP\to VP$ 是由联络 $\omega$ 确定的 $P$ 上的垂直投影。
\end{definition}

\begin{definition}
水平子丛 $HE\subset TE$ 由
\[ HE\equiv\ker\omega^V_E \]
给出。水平投影 $\omega^H_E: TE\to HE$ 是映射 $\omega^H_E\equiv\mathbf{1}_{TE}-\omega^V_E$。
\end{definition}

前推投影映射 $\dd\pi_E: TE\to TM$ 的核为 $VE$；由于 $HE$ 在 $TE$ 中横截于 $VE$，可知 $\dd\pi_E$ 在 $HE$ 上的限制是同构。

% [Tutor Note: p.145 “横截”一步是 $HE\oplus VE = TE$ 的直接推论：$\omega^H_E$ 与 $\omega^V_E$ 是互补的投影算子（$\omega^H_E+\omega^V_E=\mathbf{1}_{TE}$，$\omega^V_E\circ\omega^H_E=0$），故 $TE=HE\oplus VE$；于是 $\ker(\dd\pi_E|_{HE}) = HE\cap VE = 0$，又 $\dim HE=\dim TM=2m$，所以 $\dd\pi_E|_{HE}$ 是同构。这个同构正是下一步用 $\pi^H_\omega$ 把 $M$ 的度规拉回到 $HE$ 的依据。]

\begin{definition}
映射
\[ \pi^H_\omega \equiv \dd\pi_E|_{HE}: HE\to TM \]
是每一点处的水平同构。
\end{definition}

\begin{definition}
张量
\[
    g^H|_e\in S^2(H_eE),\qquad J^H|_e\in H_e^*E\otimes H_eE,\qquad \Omega^H|_e\in\Lambda^2(H_eE)
\]
由 $M$ 上的度规、复结构与 Kähler 形式在 $\pi^H_\omega$ 之下的拉回给出。
\end{definition}

\begin{definition}
令 $g_E$、$J_E$ 与 $\Omega_E$ 由
\begin{align*}
    g_E(v,v') &= g^V(v,v'), & & \text{for } v,v'\in VE,\\
    g_E(h,h') &= g^H(h,h'), & & \text{for } h,h'\in HE,\\
    g_E(v,h) &= 0, & & \text{for } v\in VE,\ h\in HE,
\end{align*}
定义，而 $J_E$ 与 $\Omega_E$ 类似地定义。可以证明这样定义的复结构 $J_E$ 是可积的。
\end{definition}

% [Tutor Note: p.145 “可以证明 $J_E$ 可积”是本附录唯一一处省略论证的地方（原文未给证明）。核查步骤：把 Nijenhuis 张量 $N_{J_E}(X,Y)$ 按 $X,Y$ 的 $VE$／$HE$ 分量展开，利用 $[VE,VE]\subset VE$、$[HE,VE]\subset VE$、$[HE,HE]\subset HE\oplus VE$（最后一条的 $VE$ 分量正是联络的曲率）；$(V,V)$ 与 $(H,H)$ 两块因 $J^V$、$J^H$ 各自来自真复结构 $J_X$、$J_M$ 而为零，只剩 $(H,H)$ 块给出的曲率项需要额外验证与 $J$ 相容。这也解释了为何 $J_E$ 依赖联络 $\omega$ 的选取（见 p.146 的说明）。]

总之，全空间 $E$ 被赋予 Kähler 流形的结构，其度规为 $g_E$、复结构为 $J_E$、Kähler 形式为 $\Omega_E$，均如上定义。截面及其协变导数是 $M$ 上如下整体定义的映射：

\begin{definition}
相伴纤维丛 $E=P\times_G X$ 的一个截面 $\bm{\Phi}$ 是一个映射 $\bm{\Phi}: M\to E$，它是投影映射 $\pi_E: E\to M$ 的右逆。
\end{definition}
